\documentclass[11pt]{article}

\usepackage[a4paper,margin=2.5cm]{geometry}
\usepackage{amsmath,amssymb,amsthm,mathtools}
\usepackage{enumitem}
\usepackage{hyperref}

\title{Notes on the \(h_0\)-window, \(h_0\)-eigenfunction problem}
\author{}
\date{}

\newtheorem{theorem}{Theorem}[section]
\newtheorem{proposition}[theorem]{Proposition}
\newtheorem{lemma}[theorem]{Lemma}
\newtheorem{remark}[theorem]{Remark}

\newcommand{\C}{\mathbb C}
\newcommand{\R}{\mathbb R}
\newcommand{\1}{\mathbf 1}
\newcommand{\pw}{\partial_w}
\newcommand{\pwb}{\partial_{\bar w}}

\begin{document}

\maketitle

\section{Purpose of this note}

This note records the corrected first-order reduction for the \((h_0,h_0,U)\) localization
problem and the corresponding boundary-value consequences.

The point is that, in this case, the geometry is substantially simpler than in the
\((h_1,h_1,U)\) problem. The boundary trace turns out to be constant, the interior equation is
\[
\pw u=-e^{-\pi|w|^2},
\]
and the compatibility condition on the boundary reduces directly to
\[
w\cdot \tau=0.
\]
Thus every connected \(C^1\) boundary component is a circle centered at the origin.

\section{Bargmann formulation}

Under the Bargmann transform,
\[
h_0 \longleftrightarrow e_0(w)=1.
\]
If \(T_z\) denotes the Weyl operator on Fock space, then
\[
(T_zF)(w)=e^{\pi \bar z w-\frac{\pi}{2}|z|^2}F(w-z),
\]
so
\[
T_ze_0(w)=e^{\pi \bar z w-\frac{\pi}{2}|z|^2}.
\]
Also,
\[
\langle e_0,T_ze_0\rangle_{\mathcal F^2}=e^{-\frac{\pi}{2}|z|^2}.
\]
Therefore the eigenvalue equation
\[
\mathcal L_U^{h_0}h_0=\lambda h_0
\]
is equivalent to
\begin{equation}\label{eq:bargmann-h0}
\int_U e^{\pi w\bar z}e^{-\pi|z|^2}\,dA(z)=\lambda
\qquad (w\in \C).
\end{equation}

Taking complex conjugates and replacing \(\bar w\) by \(w\), we obtain the equivalent form
\begin{equation}\label{eq:bargmann-h0-conjugated}
\int_U e^{\pi wz}e^{-\pi|z|^2}\,dA(z)=\lambda
\qquad (w\in \C).
\end{equation}

\section{The corrected compact-support equation}

Define the compactly supported distribution
\begin{equation}\label{eq:def-omega}
\omega:=
\left(-\frac{1}{\pi}\pw-\bar w\right)\!\1_U e^{-\pi|w|^2}
+\lambda\,\delta_0.
\end{equation}

\begin{proposition}\label{prop:first-order}
For every \(w\in \C\),
\[
\langle \omega,e^{\pi wz}\rangle=0.
\]
Equivalently,
\begin{equation}\label{eq:omega-expanded}
\omega
=
-e^{-\pi|w|^2}\1_U
-\frac1\pi e^{-\pi|w|^2}\pw \1_U
+\lambda\,\delta_0.
\end{equation}
\end{proposition}

\begin{proof}
Pairing against \(z\mapsto e^{\pi wz}\), one gets
\[
\left\langle
\left(-\frac{1}{\pi}\pw-\bar z\right)\!\left(e^{-\pi|z|^2}\1_U\right),
e^{\pi wz}
\right\rangle
=
\int_U ( -\bar z+w)e^{\pi wz}e^{-\pi|z|^2}\,dA(z),
\]
because \(\pw e^{\pi wz}=\pi w e^{\pi wz}\). The right-hand side equals
\[
\int_U (w-\bar z)e^{\pi wz}e^{-\pi|z|^2}\,dA(z)=0
\]
by differentiating \eqref{eq:bargmann-h0-conjugated} in \(w\). Including the \(\lambda\delta_0\)
term and using \eqref{eq:bargmann-h0-conjugated} itself yields the stated vanishing.

For the expanded form, use
\[
-\frac1\pi\pw(e^{-\pi|w|^2})-\bar w e^{-\pi|w|^2}
=
-e^{-\pi|w|^2}.
\]
\end{proof}

\begin{remark}
As in the \((h_1,h_1,U)\) note, the natural compact-support equation is first-order:
\begin{equation}\label{eq:first-order-u}
\pw u=\omega.
\end{equation}
The Paley--Wiener / Ehrenpreis--Malgrange step is adapted to \(\pw\), not directly to
\(\Delta\).
\end{remark}

\section{Boundary-value interpretation}

\begin{proposition}\label{prop:bvp}
Assume \(\partial U\) is of class \(C^1\). Let \(u\) be a compactly supported distribution
satisfying
\[
\pw u
=
-e^{-\pi|w|^2}\1_U
-\frac1\pi e^{-\pi|w|^2}\pw \1_U
+\lambda\,\delta_0.
\]
Assume moreover that \(u\) is represented by a \(C^1\)-function on \(\overline U\setminus\{0\}\),
that \(u\equiv 0\) on the exterior side of \(\partial U\), and that \(0\in U\). Then:
\begin{enumerate}[label=\textnormal{(\roman*)}]
\item the boundary trace of \(u\) is constant:
\begin{equation}\label{eq:trace}
u(w)=-\frac1\pi e^{-\pi|w|^2}
\qquad (w\in \partial U);
\end{equation}
\item in \(U\setminus\{0\}\),
\begin{equation}\label{eq:interior-first-order}
\pw u=-e^{-\pi|w|^2};
\end{equation}
hence
\begin{equation}\label{eq:interior-laplacian}
\Delta u=4\pi \bar w\,e^{-\pi|w|^2}
\qquad \text{in }U\setminus\{0\};
\end{equation}
\item on \(\partial U\),
\begin{equation}\label{eq:normal-tangent}
\partial_\nu u=-2e^{-\pi|w|^2}\nu\cdot(1,0 \text{ in complex form}),
\qquad
\partial_\tau u=-2e^{-\pi|w|^2}w\cdot \tau \ \text{after simplification};
\end{equation}
\item more importantly, differentiating the boundary trace along \(\tau\) and comparing with the
interior tangential derivative gives
\begin{equation}\label{eq:compatibility}
w\cdot \tau=0
\qquad (w\in \partial U).
\end{equation}
\end{enumerate}
\end{proposition}

\begin{proof}
Since \(u\equiv 0\) outside, one has near \(\partial U\)
\[
\pw u=\1_U\pw u+u\,\pw\1_U.
\]
Comparing coefficients of \(\pw\1_U\) gives \eqref{eq:trace}.

Away from \(\partial U\cup\{0\}\), the terms \(\pw\1_U\) and \(\delta_0\) vanish, so
\eqref{eq:interior-first-order} holds. Applying \(4\pwb\) gives \eqref{eq:interior-laplacian}.

The only geometric point needed later is the tangential derivative. Since
\[
2\pw=(\nu_1-i\nu_2)(\partial_\nu-i\partial_\tau),
\]
the interior equation \eqref{eq:interior-first-order} determines \(\partial_\tau u\) on the
boundary. On the other hand, differentiating \eqref{eq:trace} along \(\tau\) yields
\[
\partial_\tau u=2\,e^{-\pi|w|^2}(w\cdot \tau).
\]
Comparing the two expressions gives \eqref{eq:compatibility}. Since the only statement used below
is \eqref{eq:compatibility}, we omit the explicit formula for \(\partial_\nu u\).
\end{proof}

\begin{remark}
There is an even simpler way to see the compatibility. Along \(\partial U\), the trace
\eqref{eq:trace} depends only on \(|w|^2\). Hence
\[
\partial_\tau u=2e^{-\pi|w|^2}(w\cdot \tau).
\]
But the interior equation \(\pw u=-e^{-\pi|w|^2}\) forces the tangential derivative to vanish on
the boundary. Thus \(w\cdot \tau=0\).
\end{remark}

\section{Geometric consequence}

\begin{proposition}\label{prop:circle}
If \(\Gamma\subset \C\) is a connected \(C^1\) curve satisfying
\[
w\cdot \tau=0
\qquad \text{on }\Gamma,
\]
then \(|w|\) is constant on \(\Gamma\). In particular, \(\Gamma\) is contained in a circle
centered at the origin.
\end{proposition}

\begin{proof}
Let \(\gamma(s)\) be an arclength parametrization of \(\Gamma\), so \(\tau(s)=\gamma'(s)\). Then
\[
\frac{d}{ds}|\gamma(s)|^2=2\gamma(s)\cdot \tau(s).
\]
If \(w\cdot \tau=0\) on \(\Gamma\), then the right-hand side vanishes identically, hence
\(|\gamma(s)|\) is constant.
\end{proof}

\begin{theorem}\label{thm:disk}
Assume \(U\subset \C\) is bounded, simply connected, and \(\partial U\) is connected and of class
\(C^1\). Assume moreover that the corrected first-order boundary-value interpretation above is
valid. Then \(U\) is a disk centered at the origin.
\end{theorem}

\begin{proof}
By Proposition~\ref{prop:bvp}, the boundary satisfies \(w\cdot \tau=0\) everywhere. By
Proposition~\ref{prop:circle}, \(\partial U\) is contained in a circle centered at the origin.
Since \(\partial U\) is connected, it is that whole circle. Because \(U\) is simply connected, it
follows that \(U\) is the corresponding disk.
\end{proof}

\section{Final remark}

Compared with the \((h_1,h_1,U)\) problem, the corrected first-order approach is simpler here
because the boundary trace is radial with no extra polynomial factor. Once the trace and
tangential compatibility are justified, the geometry closes immediately: there is no two-branch
alternative analogous to \((1-\pi|w|^2)(w\cdot \tau)=0\).

\end{document}
